\documentclass[a4paper,11pt]{article}
\usepackage[utf8]{inputenc}
\usepackage[T1]{fontenc}
\usepackage[french]{babel}
\usepackage{lmodern}
\usepackage{amsmath,amssymb,amsthm}
\usepackage{mathrsfs}
\usepackage{enumitem}
\usepackage{multicol}

\setlist[enumerate]{label=\textbf{\textcolor{Couleur8}{\arabic*.}}, left=1em}

% Si vous voulez aussi modifier les listes enumerate imbriquées :
\setlist[enumerate,2]{label=\textcolor{Couleur8}{\alph*)}}
\setlist[enumerate,3]{label=\textcolor{Couleur8}{\roman*)}}
\setlist[enumerate,4]{label=\textcolor{Couleur8}{\Alph*)}}
\usepackage{graphicx}
\usepackage{xcolor}
\usepackage{tcolorbox}
\usepackage{geometry}
\usepackage{fancyhdr}
\usepackage{titlesec}
\usepackage{cancel}
\usepackage{ifthen}
\usepackage{tikz}
\usetikzlibrary{arrows.meta}
\usepackage{tkz-tab}
\usepackage{eurosym}

% OPTION PROF/ÉLÈVE
% Décommentez UNE des deux lignes suivantes :
\newboolean{prof}\setboolean{prof}{true}   % Version professeur (avec corrections des devoirs)
%\newboolean{prof}\setboolean{prof}{false}  % Version élève (sans corrections des devoirs)

\definecolor{Couleur1}{HTML}{4A5D7A} %Th
\definecolor{Couleur2}{HTML}{A5B5D6} %Prop
\definecolor{Couleur3}{HTML}{A8C68A} %Méthode
\definecolor{Couleur6}{HTML}{8A9CC1} %Def
\definecolor{Couleur7}{HTML}{E6B459} %Rq Voc
\definecolor{Couleur8}{HTML}{D36740} %Titres
\definecolor{correction}{HTML}{D36740} % Couleur pour les corrections
\definecolor{coopmathscorr}{HTML}{F15929} % Couleur corrections coopmaths

% Configuration des marges
\geometry{
	top=2cm,
	bottom=2cm,
	inner=1.5cm,
	outer=1.5cm,
}

% Police
\renewcommand{\familydefault}{\sfdefault}

% Compteurs
\newcounter{seancenum}
\newcounter{seancenumD}
\setcounter{seancenum}{1}
\setcounter{seancenumD}{1}

% Configuration des en-têtes et pieds de page
\pagestyle{fancy}
\fancyhf{}
\ifthenelse{\boolean{prof}}{
    \fancyhead[L]{\textcolor{Couleur8}{\textbf{Corrections Automatismes - Classe de seconde - Version Professeur}}}
}{
    \fancyhead[L]{\textcolor{Couleur8}{\textbf{Corrections Automatismes - Séance 12 - Classe de seconde}}}
}
\fancyhead[R]{\textcolor{Couleur8}{\textbf{2025-2026}}}
\fancyfoot[L]{\textcolor{Couleur8}{Mme Bertail et Mme Pinto}}
\fancyfoot[R]{\textcolor{Couleur8}{\thepage}}
\renewcommand{\headrulewidth}{1pt}
\renewcommand{\headrule}{\hbox to\headwidth{\color{Couleur8}\leaders\hrule height \headrulewidth\hfill}}

% Environnements personnalisés
\newtcolorbox{seance}[1]{
    colback=Couleur6!10,
    colframe=Couleur1!65,
    fonttitle=\bfseries\large,
    title=Correction Séance \theseancenum #1,
    left=5mm,
    right=5mm,
    top=3mm,
    bottom=3mm,
    arc=0mm,
    after=\stepcounter{seancenum}
}

\newtcolorbox{seanceD}[1]{
    colback=Couleur6!5,
    colframe=Couleur6!45,
    fonttitle=\bfseries\large,
    title=Correction Devoirs - Séance \theseancenumD #1,
    left=5mm,
    right=5mm,
    top=3mm,
    bottom=3mm,
    arc=0mm,
    after=\stepcounter{seancenumD}
}

\begin{document}

\setcounter{seancenum}{12}
\newpage
\begin{seance}{} %12
\begin{enumerate}
\item $f$ est une fonction affine, elle a donc une expression de la forme  $f(x)=ax+b$ avec $a$ et $b$ des nombres réels.
D'après le cours, on sait que pour $u\neq v$, $a=\dfrac{f(u)-f(v)}{u-v}$ Avec $u=9$ et  $v=8$, on obtient  :  $a=\dfrac{f(9)-f(8)}{9-8}=\dfrac{-27-(-25)}{9-8}=\dfrac{-2}{1}=-2$. On en déduit que la fonction $f$ s'écrit sous la forme :    $f(x)=-2 x +b.$\\On obtient $b$ en utilisant (au choix) une des deux données de l'énoncé, par exemple $f(9)=-27$.Comme $f(x)=-2x +b$, alors $f(9)=
              -2\times9+b=-18+b$ . On en déduit :\\$\begin{aligned}f(9)=-27&\iff -18+b=-27\\&\iff b=-9\\\end{aligned}$\\On en déduit $f(x)=$ ${\color[HTML]{f15929}\boldsymbol{-2x-9}}$.


\item 

\begin{enumerate}
\item On cherche les réels qui sont dans  $]11;25]$ ou bien  $[33;43]$ , ou dans les deux.\\On regarde donc la partie de l'intervalle qui est coloriée, soit en bleu, soit en rouge, soit en bleu et rouge.
Les deux ensembles sont disjoints, ils n'ont aucun élément en commun.\\
                    On ne peut pas simplifier l'écriture de $I$ qui s'écrit donc $I={\color[HTML]{f15929}\boldsymbol{]11;25]\cup[33;43]}}$.
                   \begin{tikzpicture}

    \tikzset{
      point/.style={
        thick,
        draw,
        cross out,
        inner sep=0pt,
        minimum width=2pt,
        minimum height=2pt,
      },
    }
    \clip (-2,-0.7) rectangle (15,0.5);
    	\draw[color ={black},line width = 3] (0,0)--(12,0);
	\draw[color ={red},line width = 3] (2,0)--(5,0);
	\draw[color ={blue},line width = 3] (6,0)--(9,0);
	\draw[very thick,color={red}] (1.8,0.2)--(2,0.2)--(2,-0.2)--(1.8,-0.2);
	\draw[color={red}] (2,-0.2) node[below] {$11$};
	\draw[very thick,color={red}] (4.8,0.2)--(5,0.2)--(5,-0.2)--(4.8,-0.2);
	\draw[color={red}] (5,-0.2) node[below] {$25$};
	\draw[very thick,color={blue}] (6.2,0.2)--(6,0.2)--(6,-0.2)--(6.2,-0.2);
	\draw[color={blue}] (6,-0.2) node[below] {$33$};
	\draw[very thick,color={blue}] (8.8,0.2)--(9,0.2)--(9,-0.2)--(8.8,-0.2);
	\draw[color={blue}] (9,-0.2) node[below] {$43$};

\end{tikzpicture}

\item On cherche les réels qui sont à la fois dans $[12;28]$ et dans $[14;17]$.\\On regarde la partie de l'intervalle qui est coloriée à la fois en bleu et en rouge.\\On observe que $[14;17]\subset [12;28]$ donc $I={\color[HTML]{f15929}\boldsymbol{[14;17]}}$.

\begin{tikzpicture}[baseline]

    \tikzset{
      point/.style={
        thick,
        draw,
        cross out,
        inner sep=0pt,
        minimum width=2pt,
        minimum height=2pt,
      },
    }
    \clip (-2,-0.7) rectangle (15,0.5);
    	\draw[color ={black},line width = 3] (0,0)--(12,0);
	\draw[color ={red},line width = 3] (2,-0.1)--(9,-0.1);
	\draw[color ={blue},line width = 3] (5,0.1)--(7,0.1);
	\draw[very thick,color={red}] (2.2,0.2)--(2,0.2)--(2,-0.2)--(2.2,-0.2);
	\draw[color={red}] (2,-0.2) node[below] {$12$};
	\draw[very thick,color={red}] (8.8,0.2)--(9,0.2)--(9,-0.2)--(8.8,-0.2);
	\draw[color={red}] (9,-0.2) node[below] {$28$};
	\draw[very thick,color={blue}] (5.2,0.2)--(5,0.2)--(5,-0.2)--(5.2,-0.2);
	\draw[color={blue}] (5,-0.2) node[below] {$14$};
	\draw[very thick,color={blue}] (6.8,0.2)--(7,0.2)--(7,-0.2)--(6.8,-0.2);
	\draw[color={blue}] (7,-0.2) node[below] {$17$};

\end{tikzpicture}
\end{enumerate}


\item 
 \textbf {a.}  Résolution de l'inéquation :\\
              $\begin{aligned}
              -7x-1&<-2x+3\\
              -7x-1\,{\color[HTML]{f15929}\boldsymbol{+2x}}&<-2x+3\,{\color[HTML]{f15929}\boldsymbol{+2x}}\\
              -5x-1&<3\\
              -5x&<4\\
            x&>\dfrac{4}{-5}
              \end{aligned}$
              \\
              Comme $\dfrac{4}{-5}=-\dfrac{4}{5}$, l'  ensemble $S$ des solutions de l'inéquation est
              $S= \left]-\dfrac{4}{5}\,;\,+\infty\right[$.\\\textbf {b.}  
              La courbe $\mathscr{C}_f$ est en dessous de la courbe $\mathscr{C}_g$ sur l'intervalle $\left]-\dfrac{4}{5}\,;\,+\infty\right[$.



\item  $A =  (7y-6)(-3y-5)$ \\$A =  7y\times (-3y) + 
  7y\times (-5) +
  (-6) \times (-3y)  +
  (-6) \times (-5)$ \\$A = -21y^2 + 
  (-35y) +
  18y +
  30$ \\$A =$ ${\color[HTML]{f15929}\boldsymbol{ -21y^2-17y+30 }}$


\item 
\begin{enumerate}
\item On utilise la formule du cours qui exprime le taux d'évolution $t$ en fonction de la valeur initiale $V_i$ et la valeur finale $V_f$ : $t=\dfrac{V_f-V_i}{V_i}=\dfrac{264\,600-270\,000}{270\,000}=-0{,}02=\dfrac{-2}{100}$.\\Le taux d'évolution du chiffre d'affaires est ${\color[HTML]{f15929}\boldsymbol{-2}}~\%$.

\medskip
Méthode $2$ : On arrive aussi au même résultat en passant par le coefficient multiplicateur : \\$CM=\dfrac{V_f}{V_i}=\dfrac{264\,600}{270\,000} = 0{,}98$ donc $t = 0{,}98-1= -2~\%$
\item Diminuer de $6~\%$ revient à multiplier par $1 - \dfrac{6}{100} = 1- 0{,}06 = 0{,}94$.\\Pour retrouver le prix initial, on va donc diviser le prix final par $0{,}94$. $\dfrac{922{,}14}{0{,}94}  = 981$.\\Avant la diminution le prix de mon vélo électrique était de ${\color[HTML]{f15929}\boldsymbol{981}}$\,\euro{}.
\end{enumerate}



\end{enumerate}
\end{seance}

\end{document}